

\subsection{Contrôle de concurrence multi-versions}


\begin{frame}{Algorithme de contrôle de concurrence multi-versions}

Connu sous le nom de contrôle de concurrence \textit{isolation snapshot}. Très utilisé.

\vfill

Tire parti du versionnement, qui limite les opportunités de conflits

\vfill

Très léger\,: aucun verrouillage en lecture, un contrôle sur les écritures.

\vfill

\alert{Mais}\,: ne garantit pas la sérialisabilité stricte

\end{frame}

\begin{frame}{Souvenons-nous des versions}

Un système qui fournit un niveau \texttt{repeatable read} doit s'appuyer sur un
système de versions estampillées (voir la session).

\vfill

Dans ce cas les lectures se font sur l'état de la base \alert{avant} le début de la transaction.

\vfill 

Deux possibilités de conflit entre une \alert{lecture} de $T_1$ et $T_2$.
\begin{redbullet}
  \item $r_1[d]$ est en conflit avec une écriture $w_2[d]$ 
   qui a eu lieu \alert{avant} $t_0$\,;

  \item $r_1[d]$ est en conflit avec une écriture $w_2[d]$ 
   qui a eu lieu \alert{après} $t_0$.

\end{redbullet}

\medskip

Approfondissons.  
\end{frame}

\begin{frame}{Premier cas de conflit}

Cas 1: $r_1[d]$ en conflit avec $w_2[d]$  \alert{avant} $t_0$. Alors $T_2$ a validé. Pas de risque.

\vfill

\figSlideWithSize{cc-multiversions}{13cm}

\vfill

En fait, indique que $T_2$ et $T_1$ sont en série.

\end{frame}


\begin{frame}{Second cas  de conflit}


Cas 2: $r_1[d]$  en conflit avec $w_2[d]$ 
   \alert{après} $t_0$.
   
\vfill

\figSlideWithSize{cc-multiversions2}{13cm}

\vfill
\begin{blocgauche}
 Si $T_1$ cherche
à écrire $d$ après l'écriture $w_2[d]$, un conflit
cyclique apparaît. \alert{L'algorithme doit surveiller ce cas.}
\end{blocgauche}

\end{frame}

\begin{frame}{Le contrôle de concurrence multi-versions}

\begin{defblock}{Le principe}
On vérifie, au moment d'exécuter $w_1[e]$,  qu'aucune
transaction $T_2$ n'a modifié $e$ entre le début de $T_1$ et l'instant
présent.
\end{defblock}
\vfill

En cas d'écriture $w_1[e]$, 
\begin{redbullet}
   \item si $\tau_e \leq \tau_1$ et $e$ non verrouillé\,: 
          $T_1$ verrouille (exclusif) $e$, et effectue  $w_1[e]$\,;
  \item
    si $\tau_e \leq \tau_1$ et $e$ verrouillé\,:           $T_1$ est mise en attente\,;
  \item    si $\tau_e > \tau_1$, $T_1$ est rejetée.
\end{redbullet}

\vfill

\begin{blocgauche}
\alert{Estampillage}. Au moment de $C_1$ 
   tous les tuples modifiés par $T_1$ obtiennent 
   une nouvelle version avec pour estampille l'instant du
   \texttt{commit}.
\end{blocgauche}

\end{frame}


\begin{frame}{Exemple}

Maj perdues:
$r_1(s)  r_1(c_1) r_2(s)  r_2(c_2)  w_2(s)  w_2(c_2) C_2 w_1(s) w_1(c_1) C_1$.

On prend
$\tau_1 = 100$,  $\tau_2 = 120$. 

\vfill
\begin{blocgauche}
\begin{redbullet}
  \item $T_1$ lit $s$, $T_1$ lit $c_1$,  $T_2$ lit $s$, $T_2$ lit $c_2$, sans verrouiller\,;
      \alert{aucun verrou}.
    
   \item $T_2$ veut modifier $s$\,: l'estampille de $s$
             est inférieure à  $\tau_2 = 120$: $s$ n'a pas été modifié; on 
             pose un verrou exclusif sur $s$ et on effectue
                 $w_2[s]$\,: 

   \item $T_2$ modifie $c_2$, avec pose d'un verrou exclusif\,;

   \item $T_2$ valide, relâche les verrous\,;
      versions de 
                $s$ et $c_2$ avec l'estampille $150$\,;

   \item $T_1$ veut à son tour modifier $s$\,: version de $s$ avec $\tau_s > \tau_1 = 100$. 
               $T_1$ est rejetée.
\end{redbullet}

\end{blocgauche}

\end{frame}


\begin{frame}[fragile]{Un cas de non-sérialisabilité}

L'algorithme \textit{snapshot isolation} dans certaines exécutions non sérialisables.

\vfill

Exemple\,: une simple copie d'une ligne à une autre\,: 
\begin{minted}{sql}
  function copie (id1, id2)
  begin
    val := select valeur from T where id = id1
    update T set valeur = :val where id = id2
    commit
  end
\end{minted}
\vfill 

Exécution concurrente de $copie(A,B)$ et $copie(b,A)$:
$$r_1[x]r_2[y]w_1[y]w_2[x]$$

\vfill

Non sérialisable, mais autorisée par le contrôle multi-versions.

\end{frame}


%%%%%%%%%%%
\begin{frame}{À retenir}

Algorithme très simple à mettre en œuvre, utilisé par exemple dans des contexte distribués (applications web)

\vfill

Très fluide, très peu de verrouillage, contrôle bien le cas des mises à jour perdues

\vfill

\begin{blocgauche}
Ne garantit pas la sérialisabilité stricte (des améliorations ont été proposées)
\end{blocgauche}


\end{frame}

